\documentclass[12pt]{article}
\usepackage[utf8]{inputenc}
\usepackage[T1]{fontenc}
\usepackage[margin=1in]{geometry}
\usepackage{setspace}
\usepackage{xcolor}
\usepackage[colorlinks=true, urlcolor=blue!60!black, linkcolor=blue!60!black]{hyperref}
\setlength{\parindent}{2em}
\setlength{\parskip}{1.25em}
\renewcommand{\baselinestretch}{1.0}
\title{\textit{Garcia v. San Antonio Metropolitan Transit Authority}, 469 U.S. 528 (1985)}
\author{Case Brief}
\date{}
\begin{document}
\maketitle

\section{Case Metadata}

\begin{itemize}
    \item \textbf{Case Name:} \href{https://supreme.justia.com/cases/federal/us/469/528/}{\textit{Garcia v. San Antonio Metropolitan Transit Authority}} (consolidated with No.~82-1951, \textit{Donovan, Secretary of Labor v. San Antonio Metropolitan Transit Authority et al.})
    \item \textbf{Citation:} 469 U.S. 528; 105 S. Ct. 1005; 83 L. Ed. 2d 1016 (1985)
    \item \textbf{Docket Numbers:} Nos.~82-1913 and 82-1951
    \item \textbf{Dates:} Argued March 19, 1984; Reargued October 1, 1984; Decided February 19, 1985
    \item \textbf{Court and Procedural Posture:} Supreme Court of the United States; direct appeal from the U.S. District Court for the Western District of Texas (557 F. Supp. 445 (W.D. Tex. 1983)) pursuant to 28 U.S.C.~\S~1252. Reversed and remanded.
    \item \textbf{Vote:} 5--4
    \item \textbf{Majority Opinion:} Justice Harry A. Blackmun (joined by Justices William J. Brennan, Jr., Byron R. White, Thurgood Marshall, and John Paul Stevens)
    \item \textbf{Concurring Opinions:} None
    \item \textbf{Dissenting Opinions:} Justice Lewis F. Powell, Jr. (joined by Chief Justice Warren E. Burger and Justices William H. Rehnquist and Sandra Day O'Connor) (469 U.S. at 557--579); Justice William H. Rehnquist (469 U.S. at 579--580) (also joining Powell and O'Connor); Justice Sandra Day O'Connor (joined by Justices Powell and Rehnquist) (469 U.S. at 580--589) (also joining Powell)
    \item \textbf{Constitutional and Statutory Provisions:} U.S. Const. art.~I, \S~8, cl.~3 (Commerce Clause); U.S. Const. amend.~X (Tenth Amendment); Fair Labor Standards Act of 1938 (FLSA), 29 U.S.C.~\S~201 \textit{et seq.}; Urban Mass Transportation Act of 1964 (UMTA), 49 U.S.C. App.~\S~1601 \textit{et seq.}
    \item \textbf{Prior Precedent Overruled:} \href{https://supreme.justia.com/cases/federal/us/426/833/}{\textit{National League of Cities v. Usery}}, 426 U.S. 833 (1976) (which had overruled \href{https://supreme.justia.com/cases/federal/us/392/183/}{\textit{Maryland v. Wirtz}}, 392 U.S. 183 (1968))
    \item \textbf{Primary Source Links:} \href{https://supreme.justia.com/cases/federal/us/469/528/}{Justia U.S. Supreme Court Center (Full Opinion)}; \href{https://www.law.cornell.edu/constitution-conan/amendment-10/state-sovereignty-and-tenth-amendment}{Cornell LII Constitution Annotated (State Sovereignty and the Tenth Amendment)}
\end{itemize}

\section{Facts and Procedural History}

\subsection{Factual Background and Statutory Evolution}
Passenger mass transit in San Antonio, Texas, operated privately under municipal regulation until 1959, when the city purchased the San Antonio Transit Company and established the public San Antonio Transit System (SATS). In 1978, the city transferred SATS's facilities to appellee San Antonio Metropolitan Transit Authority (SAMTA), a county-wide public transit authority that traveled over 26 million route miles and carried more than 63 million passengers between 1978 and 1980 (\href{https://supreme.justia.com/cases/federal/us/469/528/}{469 U.S. at 531}).

When SATS's finances deteriorated by 1970, it turned to federal subsidies under the Urban Mass Transportation Act of 1964 (UMTA). Between December 1970 and February 1980, SATS and SAMTA received over \$51 million in UMTA grants (more than \$31 million in capital grants and over \$20 million in operating assistance). During SAMTA's first two fiscal years, it received \$12.5 million in UMTA operating grants, \$26.8 million from local sales taxes, and \$10.1 million from fares; federal subsidies and local sales taxes together covered roughly 75 percent of operating expenses (\href{https://supreme.justia.com/cases/federal/us/469/528/}{469 U.S. at 532--533}).

Originally enacted in 1938, the Fair Labor Standards Act (FLSA) exempted state and local governments and local mass transit. Congress extended coverage to larger private transit carriers in 1961, to public schools, hospitals, and regulated mass transit in 1966 (upheld as to schools and hospitals in \href{https://supreme.justia.com/cases/federal/us/392/183/}{\textit{Maryland v. Wirtz}}, 392 U.S. 183 (1968)), and to virtually all state and local employees in 1974 while phasing out the mass transit overtime exemption. SATS complied with FLSA overtime rules until 1976, when \href{https://supreme.justia.com/cases/federal/us/426/833/}{\textit{National League of Cities v. Usery}}, 426 U.S. 833 (1976), overruled \textit{Wirtz} and barred Congress from enforcing the FLSA against States ``in areas of traditional governmental functions'' (426 U.S. at 852). Four months later, SATS notified employees it was no longer bound by FLSA overtime requirements (\href{https://supreme.justia.com/cases/federal/us/469/528/}{469 U.S. at 533--534}).

\subsection{Procedural History}
On September 17, 1979, the Department of Labor's Wage and Hour Administration issued Opinion WH-499, finding SAMTA not immune from the FLSA under \textit{National League of Cities}. On November 21, 1979, SAMTA sued the Secretary of Labor in the U.S. District Court for the Western District of Texas for declaratory relief; the Secretary counterclaimed for enforcement under 29 U.S.C.~\S~217, and SAMTA employee Joe G. Garcia intervened as a defendant (his separate overtime suit having been stayed). In December 1979, DOL formally amended 29 C.F.R.~\S~775.3(b)(3) to deny immunity to public local mass transit systems (\href{https://supreme.justia.com/cases/federal/us/469/528/}{469 U.S. at 534--535}).

On November 17, 1981, the District Court granted summary judgment to SAMTA, ruling public mass transit a traditional governmental function. While direct appeals were pending, the Supreme Court held in \href{https://supreme.justia.com/cases/federal/us/455/678/}{\textit{United Transportation Union v. Long Island R. Co.}}, 455 U.S. 678 (1982), that a state-owned commuter railroad was not a traditional governmental function, and vacated and remanded SAMTA's case (457 U.S. 1102 (1982)). On remand, the District Court adhered to its judgment (557 F. Supp. 445 (W.D. Tex. 1983)), remarking that distinguishing transit from exempt functions was like identifying pornography: ``someone knows it when they see it, but they can't describe it'' (557 F. Supp. at 453). The Supreme Court noted probable jurisdiction (464 U.S. 812 (1983)), heard argument on March 19, 1984, and ordered reargument on October 1, 1984, directing the parties to brief whether \textit{National League of Cities} should be reconsidered (468 U.S. 1213 (1984)) (\href{https://supreme.justia.com/cases/federal/us/469/528/}{469 U.S. at 535--536}).

\section{Issue(s) Presented}

\begin{enumerate}
    \item Should the Supreme Court overrule \href{https://supreme.justia.com/cases/federal/us/426/833/}{\textit{National League of Cities v. Usery}}, 426 U.S. 833 (1976), which insulated state and local governments from federal Commerce Clause regulation in ``areas of traditional governmental functions''?
    \item Does applying the FLSA's minimum-wage and overtime provisions to a publicly owned municipal mass transit authority (SAMTA) violate the Tenth Amendment or any affirmative constitutional limit on Congress's Commerce Clause power?
\end{enumerate}

\section{Holding}

\textbf{Yes} as to Issue~1; \textbf{No} as to Issue~2. In a 5--4 opinion authored by Justice Blackmun---who had provided the decisive fifth vote in \textit{National League of Cities}---the Court expressly \textbf{overruled} \href{https://supreme.justia.com/cases/federal/us/426/833/}{\textit{National League of Cities v. Usery}}. Defining state regulatory immunity by identifying ``traditional governmental functions'' is both unworkable in practice and unsound in principle. Because the Constitution protects state sovereignty primarily through procedural safeguards built into the structure of the national political process rather than judicially crafted enclaves of immunity, applying the FLSA's wage and overtime protections to SAMTA contravened no affirmative limit on Congress's Commerce Clause authority (\href{https://supreme.justia.com/cases/federal/us/469/528/}{469 U.S. at 531, 555--557}).

\section{Reasoning}

\subsection{Unworkability and Democratic Flaws of ``Traditional Governmental Functions''}
Because SAMTA conceded that its operations would easily fall within the Commerce Clause if privately owned (citing \href{https://supreme.justia.com/cases/federal/us/452/264/}{\textit{Hodel v. Virginia Surface Mining \& Reclamation Ass'n}}, 452 U.S. 264 (1981), and \href{https://supreme.justia.com/cases/federal/us/312/100/}{\textit{United States v. Darby}}, 312 U.S. 100 (1941)), immunity depended on the third \textit{Hodel} prong of \textit{National League of Cities}: whether compliance directly impaired the States' ability ``to structure integral operations in areas of traditional governmental functions'' (\href{https://supreme.justia.com/cases/federal/us/469/528/}{469 U.S. at 537--538}).

Justice Blackmun showed that eight years of lower-court litigation yielded irreconcilable results---granting immunity to ambulance regulation, driver licensing, municipal airports, solid waste disposal, and highway authorities, while denying it to traffic regulation, air transportation regulation, mental health facilities, and intrastate natural gas regulation---with no intelligible organizing principle. Analogizing to the Court's failed forty-year experiment distinguishing ``governmental'' from ``proprietary'' functions in intergovernmental tax immunity (unanimously abandoned in \textit{New York v. United States}, 326 U.S. 572 (1946)), the majority rejected both historical standards (which freeze state governance in time) and nonhistorical standards (such as ``uniquely'' or ``necessary'' services) (\href{https://supreme.justia.com/cases/federal/us/469/528/}{469 U.S. at 538--545}). More fundamentally, judicial line-drawing between ``traditional'' and ``nontraditional'' functions offends federalism itself by inviting unelected judges to penalize States when they act as Justice Brandeis's ``laboratories for social and economic experiment'' (\href{https://supreme.justia.com/cases/federal/us/469/528/}{469 U.S. at 545--547}).

\subsection{The Political Safeguards of Federalism}
Rejecting ``freestanding conceptions of state sovereignty,'' the Court turned to constitutional structure, relying on \textit{The Federalist Nos.~39, 43, 46, and 62} as well as Herbert Wechsler (\textit{The Political Safeguards of Federalism}, 54 Colum. L. Rev. 543 (1954)) and Jesse Choper (469 U.S. at 551 n.11). The Framers protected the States' residuary sovereignty primarily through the composition of the Federal Government: state control over House electoral qualifications (Art.~I, \S~2), state participation in the Electoral College (Art.~II, \S~1), and equal state representation in the Senate (Art.~I, \S~3; entrenched by Art.~V) (\href{https://supreme.justia.com/cases/federal/us/469/528/}{469 U.S. at 548--552}).

Even after the Seventeenth Amendment instituted popular election of Senators, the political process remained empirically effective: federal grants-in-aid grew from \$7 billion to \$96 billion over a quarter-century (about one-fifth of state and local expenditures), and States secured statutory exemptions from major Commerce Clause laws including the NLRA, OSHA, ERISA, the Federal Power Act, and the Sherman Act (\href{https://supreme.justia.com/cases/federal/us/469/528/}{469 U.S. at 552--554}). Consequently, the constitutional limit protecting ``States as States'' is one of \textbf{process, not result}: judicial intervention must be tailored to ``possible failings in the national political process,'' not to carving out a ``sacred province of state autonomy'' (quoting \href{https://supreme.justia.com/cases/federal/us/460/226/}{\textit{EEOC v. Wyoming}}, 460 U.S. 226, 236 (1983)). Here, SAMTA faced the same generally applicable FLSA obligations as private employers while receiving tens of millions in countervailing UMTA aid, proving the political safeguards worked as intended (\href{https://supreme.justia.com/cases/federal/us/469/528/}{469 U.S. at 554--557}).

\section{Concurring and Dissenting Opinions}

\begin{itemize}
    \item \textbf{Justice Powell, Dissenting (joined by Chief Justice Burger and Justices Rehnquist and O'Connor, \href{https://supreme.justia.com/cases/federal/us/469/528/}{469 U.S. at 557--579}):} Powell condemned the abrupt abandonment of \textit{stare decisis} after unanimous reaffirmance in \href{https://supreme.justia.com/cases/federal/us/455/678/}{\textit{Long Island}} and reaffirmance by the exact same five-Justice majority in \href{https://supreme.justia.com/cases/federal/us/456/742/}{\textit{FERC v. Mississippi}}, 456 U.S. 742 (1982). He argued that \textit{National League of Cities}---especially as framed by Blackmun's own 1976 concurrence---established a workable \textbf{balancing test}, not an \textit{a priori} taxonomy. Noting that the majority mentioned the Tenth Amendment only once, Powell charged that making Congress the sole judge of its Commerce Clause power violated \textit{Marbury v. Madison}, 5 U.S. (1 Cranch) 137 (1803), and ignored how the Seventeenth Amendment, weakened local parties, national media, and special-interest lobbying had eroded Wechsler's 1954 assumptions (469 U.S. at 565 n.9).
    \item \textbf{Justice Rehnquist, Dissenting (\href{https://supreme.justia.com/cases/federal/us/469/528/}{469 U.S. at 579--580}):} Joining both Powell and O'Connor, Rehnquist noted subtle differences among the dissenters' formulations but concluded SAMTA should prevail under any of them, declining to elaborate on ``a principle that will, I am confident, in time again command the support of a majority of this Court.''
    \item \textbf{Justice O'Connor, Dissenting (joined by Justices Powell and Rehnquist, \href{https://supreme.justia.com/cases/federal/us/469/528/}{469 U.S. at 580--589}):} Joining Powell, O'Connor grounded state autonomy in Chief Justice Marshall's Necessary and Proper Clause standard from \textit{McCulloch v. Maryland}, 17 U.S. (4 Wheat.) 316, 421 (1819): because modern Commerce Clause reach rests on the Necessary and Proper Clause, Congress's means must conform to ``the letter and spirit of the constitution,'' including the Tenth Amendment's preservation of state autonomy as part of the ``double security'' of \textit{The Federalist No.~51}. She echoed Rehnquist's prediction that the Court would eventually resume its constitutional duty.
\end{itemize}

\section{Key Quotations}

All quotations below are verified verbatim from the primary opinion text at the \href{https://supreme.justia.com/cases/federal/us/469/528/}{Justia U.S. Supreme Court Center (469 U.S. 528)}.

\begin{enumerate}
    \item \textbf{Rejection of ``Traditional Governmental Functions'' (Blackmun, J., for the Court, \href{https://supreme.justia.com/cases/federal/us/469/528/}{469 U.S. at 546--547}):}
    \begin{quote}
        ``We therefore now reject, as unsound in principle and unworkable in practice, a rule of state immunity from federal regulation that turns on a judicial appraisal of whether a particular governmental function is `integral' or `traditional.' Any such rule leads to inconsistent results at the same time that it disserves principles of democratic self-governance, and it breeds inconsistency precisely because it is divorced from those principles.''
    \end{quote}

    \item \textbf{Political Safeguards of Federalism (Blackmun, J., for the Court, \href{https://supreme.justia.com/cases/federal/us/469/528/}{469 U.S. at 550, 552}):}
    \begin{quote}
        ``Apart from the limitation on federal authority inherent in the delegated nature of Congress' Article I powers, the principal means chosen by the Framers to ensure the role of the States in the federal system lies in the structure of the Federal Government itself.'' (at 550) \dots{} ``State sovereign interests, then, are more properly protected by procedural safeguards inherent in the structure of the federal system than by judicially created limitations on federal power.'' (at 552)
    \end{quote}

    \item \textbf{Process Over Result (Blackmun, J., for the Court, \href{https://supreme.justia.com/cases/federal/us/469/528/}{469 U.S. at 554}):}
    \begin{quote}
        ``Nonetheless, against this background, we are convinced that the fundamental limitation that the constitutional scheme imposes on the Commerce Clause to protect the `States as States' is one of process, rather than one of result. Any substantive restraint on the exercise of Commerce Clause powers must find its justification in the procedural nature of this basic limitation, and it must be tailored to compensate for possible failings in the national political process, rather than to dictate a `sacred province of state autonomy.'\,''
    \end{quote}

    \item \textbf{Dissenting Critiques and Prophecy (Powell, J., at 567; Rehnquist, J., at 580; O'Connor, J., at 580, \href{https://supreme.justia.com/cases/federal/us/469/528/}{469 U.S. 528}):}
    \begin{quote}
        ``The States' role in our system of government is a matter of constitutional law, not of legislative grace.'' (Powell, J., at 567)
    \end{quote}
    \begin{quote}
        ``The Court today surveys the battle scene of federalism and sounds a retreat.'' (O'Connor, J., at 580) \dots{} ``I do not think it incumbent on those of us in dissent to spell out further the fine points of a principle that will, I am confident, in time again command the support of a majority of this Court.'' (Rehnquist, J., at 580)
    \end{quote}
\end{enumerate}

\section{Doctrinal Significance and Modern Legacy (Through 2026)}

\href{https://supreme.justia.com/cases/federal/us/469/528/}{\textit{Garcia}} remains formally unoverruled as to \textbf{generally applicable} federal regulations of state economic or employment conduct. Congress initially echoed Justice Blackmun's ``political safeguards'' thesis by enacting the Fair Labor Standards Amendments of 1985, allowing public employers to offer compensatory time off in lieu of cash overtime (see \href{https://supreme.justia.com/cases/federal/us/508/22/}{\textit{Moreau v. Klevenhagen}}, 508 U.S. 22 (1993)). The Court also reaffirmed \textit{Garcia}'s core in \href{https://supreme.justia.com/cases/federal/us/485/505/}{\textit{South Carolina v. Baker}}, 485 U.S. 505 (1988) (upholding tax rules pressuring States to issue registered bonds), and \href{https://supreme.justia.com/cases/federal/us/528/141/}{\textit{Reno v. Condon}}, 528 U.S. 141 (2000) (upholding the Driver's Privacy Protection Act because it regulated States as database suppliers rather than as sovereign lawmakers).

Yet the Rehnquist and Roberts Courts substantially fulfilled Justices Rehnquist and O'Connor's dissents---not by reviving \href{https://supreme.justia.com/cases/federal/us/426/833/}{\textit{National League of Cities}}'s ``traditional governmental functions'' test, but by encircling \textit{Garcia} with judicially enforceable federalism doctrines (\href{https://www.law.cornell.edu/constitution-conan/amendment-10/anti-commandeering-doctrine}{Cornell LII: Anti-Commandeering Doctrine}):
\begin{itemize}
    \item \textbf{Clear-Statement Rule:} \href{https://supreme.justia.com/cases/federal/us/501/452/}{\textit{Gregory v. Ashcroft}}, 501 U.S. 452 (1991), required unmistakable congressional text before interpreting a federal statute to alter core state political structures (such as mandatory retirement ages for state judges).
    \item \textbf{Anti-Commandeering:} Congress may not compel state legislatures to enact or administer federal programs (\href{https://supreme.justia.com/cases/federal/us/505/144/}{\textit{New York v. United States}}, 505 U.S. 144 (1992)), conscript state executive officers to enforce federal law (\href{https://supreme.justia.com/cases/federal/us/521/898/}{\textit{Printz v. United States}}, 521 U.S. 898 (1997)), or bar state legislatures from repealing or modifying their own laws (\href{https://supreme.justia.com/cases/federal/us/584/16-476/}{\textit{Murphy v. NCAA}}, 584 U.S. 453 (2018)). In \href{https://supreme.justia.com/cases/federal/us/599/21-376/}{\textit{Haaland v. Brackeen}}, 599 U.S. 255 (2023), the Court confirmed the boundary: evenhanded substantive standards governing both state and private actors---and ancillary state-court recordkeeping duties---do not commandeer.
    \item \textbf{Sovereign Immunity, Anti-Coercion, and Standing:} Under \href{https://supreme.justia.com/cases/federal/us/527/706/}{\textit{Alden v. Maine}}, 527 U.S. 706 (1999), Article~I cannot abrogate a nonconsenting State's sovereign immunity from private FLSA damages suits in state court, limiting enforcement to federal suits, injunctive relief, or municipal entities like SAMTA. Meanwhile, \href{https://supreme.justia.com/cases/federal/us/567/519/}{\textit{NFIB v. Sebelius}}, 567 U.S. 519 (2012), barred coercive Spending Clause conditions, and \href{https://supreme.justia.com/cases/federal/us/564/211/}{\textit{Bond v. United States}}, 564 U.S. 211 (2011), confirmed individual standing to raise Tenth Amendment challenges.
\end{itemize}

\section{Pedagogical Critiques and Exam Hypotheticals}

\subsection{Pedagogical Critiques}
\begin{itemize}
    \item \textbf{Erosion of Political Safeguards:} Echoing Justice Powell (469 U.S. at 565 n.9), critics argue that the Seventeenth Amendment, nationalized parties, and national interest groups align federal legislators with ideological coalitions rather than state institutional autonomy.
    \item \textbf{Indeterminate ``Process Failure'' Standard:} Justice Blackmun allowed judicial intervention for ``possible failings in the national political process'' (469 U.S. at 554) without defining what constitutes a process failure, risking the same subjectivity he criticized in \textit{National League of Cities}.
    \item \textbf{Formalism vs. Practical Burden:} Modern doctrine asks \textit{how} Congress regulates (general rules vs. commandeering) rather than \textit{what} function is regulated, even though generally applicable labor mandates (\textit{Garcia}) can burden state budgets far more heavily than minor executive duties (\textit{Printz}).
\end{itemize}

\subsection{Exam Hypotheticals}

\paragraph{Hypothetical 1: General Regulation vs. Commandeering.}
Congress passes a Commerce Clause statute whose Section~1 requires all public and private operators of 25+ passenger buses to electrify half their fleets by 2030, and whose Section~2 orders state DMVs to audit private operators and revoke registrations for non-compliance. \textbf{Analysis:} Section~1 survives under \textit{Garcia} and \href{https://supreme.justia.com/cases/federal/us/528/141/}{\textit{Reno}} as a generally applicable regulation of market participants (subject to \href{https://supreme.justia.com/cases/federal/us/527/706/}{\textit{Alden}} on private damages suits against States). Section~2 violates the anti-commandeering rule of \href{https://supreme.justia.com/cases/federal/us/505/144/}{\textit{New York}} and \href{https://supreme.justia.com/cases/federal/us/521/898/}{\textit{Printz}} by conscripting state executive agencies to enforce a federal regulatory scheme.

\paragraph{Hypothetical 2 (``Opposite Side'' Contrast): Federal Shift Caps for Police and Firefighters.}
Suppose Congress caps maximum shift lengths for municipal police and firefighters. \textbf{Under \textit{Garcia}:} If framed as a generally applicable safety rule covering both public and private emergency/security workers, it is valid; States must seek statutory relief through Congress (though if it targets public agencies exclusively, States may invoke the exclusivity question reserved in \textit{Reno}, 528 U.S. at 151). \textbf{Opposite Side (Under \textit{National League of Cities} / Powell \& O'Connor Dissents):} The law would be struck down regardless of general applicability, because police and fire services were paradigmatic ``traditional governmental functions'' (426 U.S. at 851), and dictating shift structures directly displaces core local fiscal and operational choices.

\end{document}
